% !TEX root = ../notes.tex

\documentclass[../notes.tex]{subfiles}

\begin{document}

\section{April 29}
Here we go.

\subsection{Whittaker Functionals for the Principal Series}
Last time, we argued for the existence of Whittaker functionals. Let's spend a moment actually constructing them.
\begin{exe}
	Fix $G\coloneqq\op{GL}_2$, and suppose that $\psi\colon F\to\CC^\times$ has conductor $\OO_F$. Choose distinct unramified characters $\chi_1,\chi_2\colon F^\times\to\CC^\times$, and let $\pi$ be the parabolic induction $V\coloneqq\op{Ind}_B^G\chi$. We construct a Whittaker functional on $V$.
\end{exe}
\begin{proof}
	One attempt to find a linear functional would be to define $\Lambda_0\colon V\to\CC$ by
	\[\Lambda_0(\varphi)\coloneqq\int_{N(F)}\varphi(u)\psi_N\left(u^{-1}\right)\,du.\]
	Certainly evaluation is a good functional on $V$, and we have taken an average to ensure that $\Lambda_0(u\varphi)=\psi_N(u)\Lambda_0(\varphi)$. However, $\varphi(u)=\varphi(1)$ for all $u$, so $\Lambda_0=0$. Instead, we could more generally take
	\[\Lambda_g(\varphi)\coloneqq\int_{N(F)}\varphi(gu)\psi_N\left(u^{-1}\right)\,du,\]
	which continues to be $N$-equivariant. If we take any $g\in N$, this vanishes for the same reason. In fact, multiplying $g$ on the left by anything in $B(F)$ will not change if $\Lambda_g$ vanishes (by the invariance of $\varphi$), and multiplying $g$ on the left by anything in $N(F)$ will also not change if $\Lambda_g$ vanishes (by the integral). Thus, we are hunting in $B(F)\backslash G(F)/N(F)$, which is $W$ by the Bruhat decomposition. This leaves us with only one non-identity element $\begin{bsmallmatrix}
		& 1 \\ 1
	\end{bsmallmatrix}$, so we are left with the candidate
	\[\Lambda(\varphi)\coloneqq\int_F\varphi\left(\begin{bmatrix}
		0 & 1 \\ 1 & 0
	\end{bmatrix}\begin{bmatrix}
		1 & x \\ 0 & 1
	\end{bmatrix}\right)\psi(-x)\,dx.\]
	This can be seen to be nonzero as follows: $\varphi$ can be chosen arbitrarily on $\begin{bsmallmatrix}
		0 & 1 \\ 1 & 0
	\end{bsmallmatrix}N(F)$, so we can arrange it for the integral to be nonzero.
\end{proof}
\begin{remark} \label{rem:whittaker-exists}
	In general, if $\dot w_0\in G$ lifts the longest element $w_0$ of $W$, then the function
	\[\varphi\mapsto\int_{F}\varphi(\dot w_0u)\psi_N\left(u^{-1}\right)\,du\]
	defines a Whittaker functional. The main issue is the convergence of this integral. The main point is that $\varphi$ failing to have compact support modulo $N(F)$ means that $\varphi$ is constant in a neighborhood of infinity, which then causes the integral to vanish due to the character $\psi_N$.
\end{remark}
\begin{remark} \label{rem:almost-uniqueness-whittaker}
	However, if $\dot w\in G$ lifts some $w\in W$ distinct from the longest element $w_0$, then one can check that the integral
	\[\Lambda_{\dot w}\varphi\mapsto\int_{F}\varphi(\dot wu)\psi_N\left(u^{-1}\right)\,du\]
	vanishes, as we would expect from \Cref{thm:whittaker-exists}. Indeed, $w\ne w_0$ means $w$ sends some simple root to another positive root. Let $\alpha_i$ be the corresponding simple root, and then we see that $\dot w\cdot u_i=\op{Ad}_wu_i\cdot\dot w$. Plugging this into the integral, we find that $\varphi(\dot wu_iu)\psi_N\left(u^{-1}\right)$ equals $\varphi(\op{Ad}_wu_i\cdot\dot wu)\psi_N\left(u^{-1}\right)$, and $\op{Ad}_wu_i$ can now be pulled out. Alternatively, we could translate the integral to $\psi_N(u_i)\varphi(\dot wu_iu)\psi_N\left(u^{-1}\right)$, so we see that
	\[\Lambda_{\dot w}(\varphi)=\psi_N(u_i)\Lambda_{\dot w}(\varphi),\]
	so $\Lambda_{\dot w}(\varphi)=0$ follows.
\end{remark}
We have thus sketched the existence of Whittaker functionals on these principal series, and we have a suggestion of uniqueness because none of the other integrals we could try would give us a new Whittaker functional. It is possible to upgrade this into a proof of uniqueness.
\begin{proposition}
	Fix a pinned reductive group $G$ over a nonarchimedean local field $F$. Then $\op{Ind}_B^G\chi$ admits a nonzero Whittaker functional, which is unique up to scalar.
\end{proposition}
\begin{proof}
	Existence follows from \Cref{rem:whittaker-exists}, so we focus on uniqueness. Choose a Whittaker functional $\Lambda$ of $\op{Ind}_B^G(\chi)$. Note that there is a projection from $\mc H_G\coloneqq C_c^\infty(G(F))$ to $\op{Ind}_B^G\chi$ given by averaging over $B$ by $\chi\delta^{1/2}$. Thus, the composite
	\[\mc H_G\to\op{Ind}_B^G\chi\stackrel\Lambda\to\CC\]
	provides a distribution $D$ for which $D(L(b)R(u)f)=(\chi\delta^{1/2})(b)\psi_N(u)D(f)$ for any $b\in B(F)$ and $u\in N(F)$. It now turns out that the space of such invariant distributions on $G(F)$ is one-dimensional and given by
	\[f\mapsto\int_{B(F)\times N(F)}f(b\dot w_0u)\left(\chi\delta^{1/2}\right)(b)\psi_N(u)\,db\,du.\]
	This completes the proof.
\end{proof}
\begin{remark}
	The same idea as in \Cref{rem:almost-uniqueness-whittaker} shows that $\op{Ind}_P^G(\chi)$ admits no Whittaker functional when $\chi$ is one-dimensional and $P$ is some parabolic subgroup strictly larger than $B$.
\end{remark}

\subsection{Uniqueness of Whittaker Functionals}
We will now show that Whittaker models are unique for nonarchimedean local fields. Over finite fields (\Cref{thm:gelfand-graev}), the idea was to show that the endomorphism algebra, presented as a convolution algebra, is commutative, which was done by finding a suitable anti-involution on the group.

However, to make sense of a convolution algebra, we should work with distributions. Throughout, we fix a split reductive group $G$ over a local field $F$, which we equip with a Borel subgroup $B\subseteq G$ with a unipotent radical $N$.
\begin{notation}
	Fix everything as above. Let $\mc D^{N(F)\times N(F)}$ denote the distributions on $G(F)$ which are eigenvectors for the left- and right-translation action by $N(F)\times N(F)$ with eigenvalue $\psi_N\boxtimes\psi_N^{-1}$.
\end{notation}
\begin{remark}
	The inverse is needed to ensure that right translation is given by a left action.
\end{remark}
\begin{remark}
	If a distribution $D$ comes from a function $\Phi$, then
	\[\Phi(u_1gu_2)=\psi_N(u_1)^{-1}\psi_N(u_2)^{-1}\Phi(g).\]
\end{remark}
Now, distributions are not an algebra (indeed, they should be a co-algebra), but we can still hunt for an anti-involution $\iota$ on $G(F)$, which acts on our distributions by
\[(\iota D)(f)\coloneqq D(f\circ\iota).\]
Thus, we have an involution on our distribution space. As before, we will take $\iota$ to be $\dot w_0\tau(g)^{-1}\dot w_0^{-1}$, where $w_0$ is the longest Weyl element, $\dot w_0$ is a lift, and $\tau$ is the Cartan involution.
\begin{proposition} \label{prop:distribution-fixed}
	Fix everything as above. Then the action of $\iota$ on $\mc D^{N(F)\times N(F)}$ is the identity.
\end{proposition}
\begin{proof}[Idea]
	The moral is that $\iota$ fixes all ``relevant'' double cosets, exactly as was done for finite fields.
\end{proof}
We are now ready to prove uniqueness.
\begin{proof}[Proof of uniqueness in \Cref{thm:whittaker-exists}]
	Fix a smooth irreducible representation $\pi\colon G(F)\to\op{GL}(V)$.
	\begin{enumerate}
		\item Suppose that we are given Whittaker functionals $\Lambda\colon V\to\CC$ and $\Lambda'\colon V^\lor\to\CC$. For motivation, note that if $\Lambda$ and $\Lambda'$ were in fact smooth, then the function $g\mapsto\langle g\Lambda,\Lambda'\rangle$ would be a function which is an eigenvector for $N(F)$ on both the left and right on $\psi_N$, so this function is invariant under $\iota$.

		In general, $\Lambda$ and $\Lambda'$ do not have to be smooth, but we can still recover a distribution as
		\[\mu(\varphi)\coloneqq\langle\Lambda*\varphi,\Lambda'\rangle,\]
		for $\varphi\in\mc H_G$. Here, $\Lambda*\varphi$ is the functional $v\mapsto\Lambda(\varphi\cdot v)$, where $\varphi\cdot v$. Because $\varphi$ is smooth, it turns out that the convolution $\Lambda*\varphi$ is now smooth: $\varphi$ being right-invariant by a compact open subgroup gives the same for $\Lambda*\varphi$. Thus, we can pair with $\Lambda'$, as needed.

		This $\mu$ is basically a matrix coefficient, and one can check that the eigenvector properties for $\Lambda$ and $\Lambda'$ give the corresponding properties for $\mu$, so $\mu\in\mc D^{N(F)\times N(F)}$.

		\item We would now like to build a matrix coefficient given two nonzero Whittaker functionals $\Lambda_1,\Lambda_2\colon V\to\CC$. To this end, we must identify $V^\lor$ and $V$ in some way, which is the contragredient: one can define $\pi^*\colon G(F)\to\op{GL}(V)$ by $\pi^*(g)\coloneqq\pi\left(\iota g^{-1}\right)$. We now claim that $\pi^*$ is isomorphic to $V^\lor$. Indeed, for finite-dimensional representation, this can be seen on the level of the weight decomposition, where the point is that the map $g\mapsto\dot w_0\tau(g)^{-1}\dot w_0^{-1}$ sends a dominant integral weight to its dual dominant integral weight because $\tau$ acts by $-w_0$ on the weights.
		
		In the general case, we should compare characters: for any $\varphi\in C_c^\infty(K\backslash G(F)/K)$, we can define $\tr\pi(\varphi)$ as the trace on $V^K$, so $\tr\pi$ upgrades to a distribution of $G(F)$, which we will denote $\Theta_\pi$. Now, the characters of $\pi^\lor$ and $\pi^*$ can both be related to the character of $\pi$: on one hand, one has
		\[\Theta_{\pi^\lor}(f)=\Theta_\pi\left(f\circ(-)^{-1}\right)\]
		directly from the definition by passing to the finite-dimensional representations. On the other hand, the definition of $\pi^*$ gives
		\[\Theta_{\pi^*}(f)=\Theta_\pi\left(f\circ(-)^{-1}\circ\iota\right),\]
		so it only remains to check that composing with $\iota$ does not change the character. However, $\iota$ fixes every conjugacy class.

		\item Thus, we may define $\Lambda_2'\colon V^\lor\to\CC$ by using the identification of $V^\lor$ with the representation $\pi^*$. Then $\Lambda_2'$ is a Whittaker functional for $\psi_N^{-1}$. It follows that the corresponding matrix coefficient $\mu_{12}(\varphi)\coloneqq\langle\Lambda_1*\varphi,\Lambda_2'\rangle$ is successfully in $\mc D^{N(F)\times N(F)}$. One can similarly define $\mu_{21}$ by the symmetric definition, and one can check that $\mu_{21}$ and $\mu_{12}$ satisfy
		\[\mu_{21}=\iota\mu_{12}.\]
		Thus, \Cref{prop:distribution-fixed} tells us that $\mu_{12}=\mu_{21}$, so
		\[\langle\Lambda_1*\varphi,\Lambda_2'\rangle=\langle\Lambda_2*\varphi,\Lambda_1'\rangle.\]

		\item We complete the proof. We claim that the two functions $(\Lambda_1*-),(\Lambda_2*-)\colon\mc H_G\to V^\lor$ have the same kernel. It then follows that both maps factor as isomorphisms through the same quotient of $\mc H_G$, which completes the proof because isomorphisms are unique up to scalar by Schur's lemma.

		By symmetry, it is enough to just show one inclusion, so suppose that $\Lambda_1*\varphi=0$ implies that $\Lambda_2*\varphi=0$. Well, if $\Lambda_1*\varphi=0$, we find that $\Lambda_1*\varphi*\varphi'=0$ for any $\varphi'$, so
		\[\langle\Lambda_1*\varphi*\varphi',\Lambda_2'\rangle=0.\]
		From the previous step, we conclude that
		\[\langle\Lambda_2*\varphi*\varphi',\Lambda_1'\rangle=0.\]
		We would now like to check that $\Lambda_2*\varphi=0$. Well, if $\Lambda_2*\varphi$ is nonzero, then the vectors $\Lambda_2*\varphi*\varphi'$ span $V^\lor$, so pairing with $\Lambda_1'$ must get a nonzero value somewhere, which is a contradiction!
		\qedhere
	\end{enumerate}
\end{proof}

\end{document}